\section{Proof of the ordinary-reduction theorem}\label{sec:ordinarity}

Fix a complex embedding of the base field. The rational Betti \(H^1\) is
two-dimensional over \(E\), and Hodge conjugation exchanges the embeddings
in each pair. The two holomorphic multiplicities of such a pair therefore
sum to two. We call the pair \emph{balanced} when both are one.

\begin{proof}[Proof of Theorem~\ref{lem:rank-two}]
Work over a finite extension containing a normal closure of \(E\), with the
endomorphisms and polarization defined. Choose a rational prime \(\ell\)
away from the discriminant of \(E\) and the index of the endomorphism
order in \(\mathcal O_E\). Take \(N\) with \(\ell^N>2^{2e+1}\) and further
extend the base to make \(J[\ell^N]\) rational and to contain the
\(\ell^N\)-th roots of unity. The absolute degree-one primes have density one,
since primes of higher residue degree have norm at least the square of
their underlying rational prime.
At a good such prime of residue characteristic \(p\), outside a finite
set, \(p\ne\ell\) and \(E\otimes\mathbb Q_p\) splits. The actual
endomorphism order in \(E\) has completion \(\mathcal O_E\otimes\mathbb Z_p\)
away from finitely many primes. Its idempotents split the \(p\)-divisible
group, and the crystalline summands have rank two.

Here are the trace comparisons used below. For an actual endomorphism
\(h\) of the reduction, the characteristic polynomial on \(H^1\) is
the degree polynomial of \(h\), both in \(\ell\)-adic and in crystalline
cohomology. Its value at a sufficiently large integer \(m\) equals
\(\deg([m]-h)\), by the exterior algebra description of cohomology
and the action on top cohomology. Take a rational basis of \(E\) made
of integral endomorphisms \(a\) and apply this to \(a\pi_p\), where
\(\pi_p\) is the \(p\)-power Frobenius of the reduction. The
nondegenerate field-trace pairing identifies the relative traces in
the two cohomologies as the scalar extensions of one element
\(\tau\in E\), determined by
\(\Tr_{E/\mathbb Q}(a\tau)=\Tr(a\pi_p\mid H^1)\).
The same argument with \(a\) alone identifies the ranks of the
crystalline summands with the rank-two Betti comparison ranks.
Every embedding of \(\tau\) gives a sum of two Frobenius eigenvalues,
which are algebraic integers of absolute value \(\sqrt p\) under
every complex embedding. Hence \(\tau\in\mathcal O_E\) and
\[
                             |N_{E/\mathbb Q}(\tau)|\leq 2^{2e}p^e.
\]

The split crystalline components have slopes in \([0,1]\), with sum
given by the dimension of the corresponding \(p\)-divisible factor.
The dimensions agree with the characteristic-zero Hodge
multiplicities in the split idempotent pieces, since the ranks in
the Lie algebra are preserved by good reduction. In each conjugate pair
of type \((0,2)\), the component of dimension two has both slopes one.
Its trace is divisible by \(p\).
In a nonordinary pair of type \((1,1)\), the two slopes of the
nonordinary component lie strictly between zero and one, since they
sum to one. The polarization pairs that component with slopes obtained
by subtracting its slopes from one. Both component traces are therefore
divisible by \(p\). The slopes show that each trace has positive
\(p\)-adic valuation, and integrality of the relative trace at each
split place makes that valuation at least one.
If \(b\) is the number of pairs of type \((1,1)\), nonordinarity implies
\[
                         p^{e-b+2}\mid N_{E/\mathbb Q}(\tau).
\]
Indeed all \(e-b+2\) places counted here are distinct and have degree one.
For \(b\leq2\) the normalized trace norm belongs to a fixed finite set
of integers at every nonordinary degree-one prime.

We now use a torsion congruence, as in the ordinary-reduction argument
attributed to Katz by Ogus \cite[Proposition~2.7.2]{ogus}. At our good
degree-one prime, Frobenius acts as the identity modulo \(\ell^N\) on
the Tate module, and \(p\equiv1\pmod{\ell^N}\). For every element \(a\)
of the integral endomorphism order in \(E\) we obtain
\[
  \Tr_{E/\mathbb Q}(a\tau)
    =\Tr(a\pi_p\mid H^1)
    \equiv\Tr(a\mid H^1)
    =\Tr_{E/\mathbb Q}(2a)
    \pmod{\ell^N}.
\]
The order completes at \(\ell\) to \(\mathcal O_E\otimes\mathbb Z_\ell\),
where the field-trace pairing is perfect by the choice of \(\ell\).
Thus \(\tau\equiv2\) in this completion modulo \(\ell^N\). Taking the
norm and using the congruence for \(p\) gives
\[
 p^{-e}N_{E/\mathbb Q}(\tau) \equiv 2^{2e}\pmod{\ell^N}.
\]
If reduction were nonordinary, the integer on the left would lie in
\([-2^{2e},2^{2e}]\). The modulus forces it to equal \(2^{2e}\).
Equality in the product Weil bound forces \(|\sigma(\tau)|=2\sqrt p\)
for every complex embedding \(\sigma\) of \(E\), so
\(\tau\overline{\tau}=4p\) in \(E\). At either split place of a
nonordinary balanced pair, the left side has valuation at least two
because both traces are divisible by \(p\). The right side has
valuation one for \(p\ne2\). This contradiction excludes all
nonordinary absolute degree-one primes outside the finite exclusions
over the chosen extension, proving the density claim.
\end{proof}

\subsection{Comparison with ordinary-reduction theorems}

Sawin determines the density of ordinary primes for abelian surfaces and
uses a Weil bound and algebraic Chebotarev to control an integral,
normalized Frobenius trace \cite[Theorem~1 and Corollary~2]{sawin}.
Fit\'e's results apply to certain threefolds, fourfolds, and higher-dimensional
varieties with extra endomorphisms. In particular, a Galois endomorphism
subfield of degree \(g=\dim J\) gives positive density of \(p\)-rank
at least \(\lceil g/2\rceil\) \cite[Theorems~1--4 and Corollary~5]{fite}.
These results supply several small-dimensional cases of the cover
argument.

Cantoral Farf\'an, Li, Mantovan, Pries, and Tang use a relative trace norm
to prove density-one ordinarity after finite extension for simple
signature with full geometric endomorphism algebra a CM field
\cite[Assumptions~2.1 and~2.5, Proposition~3.12, and Theorem~4.18]{cantoral}.
In relative rank two, simple signature permits one balanced conjugate
pair. Theorem~\ref{lem:rank-two} extends the potentially density-one part
of that rank-two result to two balanced pairs and to an acting field
\(E\) inside an arbitrary full geometric endomorphism algebra. The
additional divisibility is the trace contribution of the second
member of the nonordinary balanced pair. The proof works at split
primes after base extension. The results of \cite[Theorems~4.18--4.19]{cantoral}
also address higher relative rank, ordinary density over the original
field in certain cases, and \(\mu\)-ordinary reduction.

The primitive factors of cyclic orders \(7,9,11,15\) in the quintic
argument have dimensions \(6,6,10,8\), respectively, and the indicated
cyclotomic action has exactly two balanced pairs in each case. The
theorem applies to each smooth member with those branch exponents,
including members with additional endomorphisms. By comparison, the
order-five four-branch case has one balanced pair. Its exponents
\((4,2,2,2)\) become \((1,3,3,3)\) on changing the deck generator,
giving the Moonen family treated in
\cite[Corollary~5.1 and \S5.2]{cantoral} under the endomorphism hypotheses
there.
